% Publication-ready G18 continuation.  This file is input by
% G18_FIXED_SPACING_CARRIER_BRIDGE_INSERT.tex and assumes its repaired CMP(1)
% close-projection/GNS transport lemma and complete fixed-spacing
% three-component band theorem.
%
% Bibliography addition required if this insert is merged:
% \bibitem{Yarotsky2006GroundStates} D. A. Yarotsky,
%   \emph{Ground states in relatively bounded quantum perturbations of
%   classical lattice systems}, Commun. Math. Phys. \textbf{261}, 799--819
%   (2006), \href{https://arxiv.org/abs/math-ph/0412040}
%   {arXiv:math-ph/0412040}.

\subsection{Exact isolation of the internal carrier sheet away from
\texorpdfstring{$\Gamma$}{Gamma}}
\label{sec:g18-internal-sheet-closure}

This subsection works at fixed lattice spacing in the $SU(3)$ Hamiltonian
theory and in the unrescaled energy convention of Eq.~\eqref{eq:H}.  It uses
the repaired close-projection theorem and
Theorem~\ref{thm:g18-fixed-spacing-band}; no Wilson-action transfer matrix is
used.  The result closes the internal rank-one sheet on every fixed momentum
patch whose closure avoids $\Gamma$.  It does not give a gap uniform as
$k\to0$.

For $\mu>0$, let $\mathcal A_\mu$ denote the unital Banach $*$-algebra of
$3\times3$ convolution kernels $K=(K_{\alpha\beta}(x))$ on $\mathbb Z^3$ with
norm
\begin{equation}
 \|K\|_{\mu,\sharp}:=
 \max\left\{
 \max_\alpha\sum_{x,\beta}e^{\mu|x|_1}|K_{\alpha\beta}(x)|,
 \max_\beta\sum_{x,\alpha}e^{\mu|x|_1}|K_{\alpha\beta}(x)|
 \right\}.
 \label{eq:g18-weighted-algebra}
\end{equation}
Its Fourier transform is analytic for
$|\operatorname{Im}k_j|<\mu$ and
$\sup_{k\in\mathbb T^3}\|K(k)\|\le\|K\|_{\mu,\sharp}$.

\begin{lemma}[Coupling-analytic exponentially local symbol]
\label{lem:g18-u-analytic-symbol}
There are $\rho_0>0$ and $\mu>0$ such that the orthonormal-frame Hamiltonian
of Theorem~\ref{thm:g18-fixed-spacing-band} is represented by a map
\begin{equation}
 u\longmapsto h_u\in\mathcal A_\mu
 \label{eq:g18-banach-holomorphic}
\end{equation}
which is holomorphic for complex $|u|<\rho_0$.  For real $u$, $h_u(k)$ is
Hermitian.  The same statement holds for the canonical Gram kernel $G_u$ and
the physical one-mark energy kernel $\mathcal K_u$, and
\begin{equation}
 h_u=G_u^{-1/2}\mathcal K_uG_u^{-1/2}.
 \label{eq:g18-h-from-gk}
\end{equation}
Here the inverse square root is the branch satisfying $G_0^{-1/2}=I$.
\end{lemma}

\begin{proof}
Choose $\rho_0$ strictly inside both the repaired CMP(1) smallness domain and
the complex ground-state contraction domain.  With
$\lambda_0=27c_2\rho_0<1/33$, one may use the fixed contour
\[
 \gamma_{\rho_0}:\quad |z-4|=(1-\lambda_0)/8
\]
in the gap-one rescaling.  It encloses the target cluster for every
$|u|\le\rho_0$ and stays in the coefficient resolvent set.  Indeed, the
target disc has radius at most $4\lambda_0$, while the nearest other free
centres are $7/2$ and $17/4$.  The inequalities separating all three discs
from this contour reduce on the upper side to
$33\lambda_0<1$ (and on the lower side to the weaker
$\lambda_0<1/9$).  For complex $u$ this statement uses the coefficient
Neumann-resolvent bound, whose absolute majorant depends on $|u|$, rather
than a self-adjoint spectral-ordering argument.  This fixed contour, rather
than a contour depending on the evaluation point $u$, is what makes the
following holomorphic calculation literal.

The product-ground-state fixed-point map used in
\cite{Yarotsky2004QP} is analytic in the bounded interaction.  On a closed
complex $u$-disc contained in the contraction domain its contraction
constant and all coefficient majorants depend only on $|u|$.  The analytic
contraction theorem therefore makes its fixed point, the dressed coefficient
perturbation $F_u$, and the Neumann resolvent series holomorphic in their
weighted coefficient norms.  Contour integration on $\gamma_{\rho_0}$ makes
the coefficient Riesz projection $\mathcal P_u^-$ holomorphic in every
slightly weaker rooted norm.

We record why the spatial estimate remains valid on a complex disc.  In the
proof of \cite[Theorem~1]{Yarotsky2006GroundStates}, Lemmas~3--4 give
holomorphic polymer activities whose absolute values, uniformly on a closed
parameter disc, obey a majorant of the form
\[
 |a_u(\chi)|\le C_\rho\varepsilon_\rho^{|\operatorname{supp}\chi|},
 \qquad \varepsilon_\rho<1,
\]
after the coupling disc is made small enough.  The number of polymers of
size $n$ through a specified point grows at most exponentially in $n$.
Consequently the connected-cluster series is absolutely and locally
uniformly convergent for complex $u$, which is the argument used there to
prove weak-star analyticity.  With two marked local insertions, every
connected cluster joining supports $I$ and $J$ has size bounded below by a
constant times $\operatorname{dist}(I,J)$.  Summing the same activity
majorant therefore gives the complex-disc estimate
\begin{equation}
 \sup_{|u|\le\rho}
 |\omega_u^{\mathbb C}(A^*B)
   -\omega_u^{\mathbb C}(A^*)\omega_u^{\mathbb C}(B)|
 \le C_\rho^{|I|+|J|}\tau_\rho^{\operatorname{dist}(I,J)}
       \|A\|\|B\|,
 \qquad \tau_\rho<1.
 \label{eq:g18-complex-cluster-majorant}
\end{equation}
Here $\omega_u^{\mathbb C}$ denotes the holomorphic continuation; for real
$u$ it is the physical ground state.  This is an absolute cluster bound, not
an inference from real-$u$ positivity or from pointwise weak-star
analyticity.

Put $\mathbb D_u:=D+F_u$ for the gap-one coefficient Hamiltonian
and use the one-mark identities
\begin{align}
 (G_u)_{st}
 &=\omega_u\!\left(\widehat w_s^*\,\widehat{\mathcal P_u^-w_t}\right),
 \label{eq:g18-one-mark-g-analytic}\\
 (\widehat{\mathcal K}_u)_{st}
 &=\omega_u\!\left(\widehat w_s^*\,
       \widehat{\mathbb D_u\mathcal P_u^-w_t}\right),
 \qquad \mathcal K_u:=\frac23\widehat{\mathcal K}_u.
 \label{eq:g18-one-mark-k-analytic}
\end{align}
They agree with the Hilbert-space Gram and energy kernels for real $u$ but,
unlike a two-mark norm formula, also give their holomorphic continuation to
complex $u$.  To make the summation explicit, write
$\mathcal P_u^-w_t=w_t+c_t(u)$.  The rooted Neumann majorant gives, uniformly
for $|u|\le\rho<\rho_0$,
\[
 \sum_J\eta^{-d(J,S_t)}\|D_Jc_{t,J}(u)\|\le A_{\eta,\rho},
\]
and the same calculation applied to
\[
 \mathbb D_u\mathcal P_u^-w_t
 =4w_t+Dc_t+F_uw_t+F_uc_t
\]
gives a rooted majorant $A'_{\eta,\rho}$.  Both coefficient families are
holomorphic.  Exact-support projections commute with charge conjugation, so
$w_t$, $c_{t,J}$ and every displayed energy correction are charge odd.  The
analytically continued ground-state functional is charge even; hence both
one-point terms in
\eqref{eq:g18-complex-cluster-majorant} vanish identically, for complex as
well as real $u$.  For separated roots apply
\eqref{eq:g18-complex-cluster-majorant} term by term and use
$\operatorname{dist}(S_s,S_t)\le
\operatorname{dist}(S_s,J)+d(J,S_t)$, exactly as in the real CMP(3) Schur
sum.  Contacting roots form a fixed finite list and are bounded directly by
the absolutely convergent marked-cluster series.  Choose $\eta$ and then the
disc so that
$\theta_\rho:=\max\{C_\rho\eta,\tau_\rho\}<1$.  Since the number of roots at
distance $n$ grows quadratically, any
$0<\mu_0<-\log\theta_\rho$ gives constants $C,\mu_0>0$ such that
\[
 \sum_{x\in\mathbb Z^3}e^{\mu_0|x|_1}
 \sup_{|u|\le\rho}\bigl(
   |G_{\alpha\beta}(x;u)|
   +|\mathcal K_{\alpha\beta}(x;u)|
 \bigr)\le C,
 \qquad \rho<\rho_0.
\]
Finite spatial truncations are holomorphic, and this summable supremum
majorant makes them converge locally uniformly in $\mathcal A_\mu$ for any
$0<\mu<\mu_0$.  Thus $G_u$ and $\mathcal K_u$ are Banach-valued holomorphic
maps.

After shrinking $\rho_0$ if necessary,
$\|G_u-I\|_{\mu,\sharp}<1$.  The binomial series for $G_u^{-1/2}$ converges
locally uniformly in the unital Banach algebra $\mathcal A_\mu$ and is
holomorphic.  Equation~\eqref{eq:g18-h-from-gk} proves the claim.  The
coefficient calculation in the gap-one Hamiltonian
$\widehat H=(3/2)H$ would give the non-scalar coefficient
$(3/2)t_3=5/408$.  Multiplying the completed symbol by $2/3$ returns exactly
$t_3=5/612$ in the unrescaled normalization of Eq.~\eqref{eq:H}.
\end{proof}

\begin{lemma}[Finite-to-infinite second-order coefficient]
\label{lem:g18-finite-infinite-u2}
In the orthonormal frame of
Theorem~\ref{thm:g18-fixed-spacing-band},
\begin{equation}
 [u^0]h_u(k)=\frac83I_3,\qquad
 [u^1]h_u(k)=I_3,\qquad
 [u^2]h_u(k)=\sigma_2I_3+\frac5{612}B(k)B(k)^\dagger,
 \label{eq:g18-exact-infinite-u2}
\end{equation}
where $\sigma_2$ is a real scalar.  If the separate $SU(3)$ diagonal ledger
quoted in Section~\ref{sec:su3} is adopted, then $\sigma_2=11/306$; the
internal-sheet conclusion below does not use this value.
\end{lemma}

\begin{proof}
First work in a sufficiently large finite periodic box ($L\ge5$), in the
physical charge-odd trivial-flux shell, and use a fixed contour valid on a
complex $u$-disc.  Kato parallel transport gives a holomorphic intertwiner
from the free plaquette shell to the interacting Riesz island; on the real
axis it is unitary.  Differentiating its reduced operator at $u=0$ gives the ordinary
degenerate reduced-resolvent coefficient
\[
 P_-(-V)Q_-(E_{F,3}-H_E)^{-1}Q_-(-V)P_-
   -[u^2]E_{\rm vac,L}(u)I,
\]
with the conventional scalar terms absorbed into the diagonal ledger.  This
is precisely the $Q_-=1_--P_-$ prescription of Eq.~\eqref{eq:heff}, rather
than an eigenvalue-only comparison.  The retained-shell theorem,
the first-order charge/determinant census, and
Corollary~\ref{cor:assembly} therefore give
\[
 k_{L,0}=\frac83I,\qquad k_{L,1}=I,\qquad
 k_{L,2}=\sigma_{2,L}I+\frac5{612}BB^\dagger.
 \label{eq:g18-finite-kato-u2}
\]

The canonical Gram orthonormalization used to define $h_{L,u}$ is another
holomorphic intertwiner which equals the free plaquette frame at $u=0$ and
is unitary for real $u$.  The two identifications differ by a holomorphic
near-identity similarity $Z(u)$, unitary on the real axis.  If
$Z_0=I$ and
\[
 a(u)=a_0I+ua_1I+u^2a_2+\sum_{n\ge3}u^na_n,
\]
then direct expansion gives
\[
 [u^2]\,Z(u)^{-1}a(u)Z(u)=a_2;
\]
all possible frame terms are commutators with $a_0I$ or $a_1I$.  Thus the
full second-order matrix, not merely its eigenvalues, is unchanged by the
orthonormal-frame choice.  This is the point at which scalarity of both the
zeroth- and first-order coefficients is essential.

It remains to justify passage to infinite volume without confusing torus
convolution with convolution on $\mathbb Z^3$.  Let $\mathcal A_{\mu,L}$ be
the corresponding weighted convolution algebra on $\mathbb Z_L^3$, using
the torus distance, and let
\[
 (\Pi_LK)([x])=\sum_{n\in\mathbb Z^3}K(x+nL)
\]
be periodisation.  Then $\Pi_L:\mathcal A_\mu\to\mathcal A_{\mu,L}$ is a
contractive algebra homomorphism.  Let $G_{L,u}$ and
$\mathcal K_{L,u}$ be the periodic finite-volume one-mark kernels.  The
coefficient fixed-point and marked-cluster series converge term by term to
their infinite-volume series for every fixed centred displacement, locally
uniformly on a smaller complex $u$-disc.  The periodic ground-state limit is
the same state as the open-box limit because the small-perturbation ground
state is unique \cite{Yarotsky2005Uniqueness}.  This first identifies the
two limits for real $u$; local uniform holomorphy and the identity theorem
then identify their complex continuations on the common disc.

The proof of Lemma~\ref{lem:g18-u-analytic-symbol} gives the stronger bound,
uniformly in $L$,
\[
 \sum_{[x]}e^{\mu_0d_L([x],0)}
 \sup_{|u|\le\rho}
 \bigl(|G_{L,\alpha\beta}(x;u)|
       +|\mathcal K_{L,\alpha\beta}(x;u)|\bigr)\le C_\rho.
\]
Fix $0<\mu<\mu_0$.  Split the torus norm at
$d_L([x],0)=R$, with $2R<L$.  The finite part of
$G_{L,u}-\Pi_LG_u$ (and of the energy kernel) converges locally uniformly by
the preceding termwise limit.  Both tails, including every periodised image,
are at most $C_\rho e^{-(\mu_0-\mu)R}$.  Letting first $L\to\infty$ and then
$R\to\infty$ proves
\[
 \|G_{L,u}-\Pi_LG_u\|_{\mu,\sharp,L}\longrightarrow0,\qquad
 \|\mathcal K_{L,u}-\Pi_L\mathcal K_u\|_{\mu,\sharp,L}
 \longrightarrow0
\]
locally uniformly in $u$.  On a smaller disc all Gram kernels obey a common
$\|G_{L,u}-I\|_{\mu,\sharp,L}\le r_G<1$.  The binomial series is therefore
uniform in $L$, and the homomorphism property of $\Pi_L$ gives
\[
 \|h_{L,u}-\Pi_Lh_u\|_{\mu,\sharp,L}\longrightarrow0
\]
locally uniformly.  This also shows explicitly that wrap-around products in
the finite-torus inverse square root are accounted for by periodisation, not
silently identified with infinite convolution.  Cauchy's formula on any
circle inside that disc now gives
\begin{equation}
 [u^2]h_u(x) =
 \lim_{L\to\infty}[u^2]h_{L,u}(x)
 \label{eq:g18-cauchy-coefficient-limit}
\end{equation}
for each fixed displacement $x$.  This is stronger than convergence at a
few real couplings: local uniformity on a complex contour is what permits
interchanging the thermodynamic limit with Taylor-coefficient extraction.

At order two every connected process is supported on one face or on two
faces sharing one link.  Hence, once the roots and this finite process
neighbourhood fit without wrapping, the finite-volume coefficient is exactly
the infinite-lattice coefficient.  Lemma~\ref{lem:process-complete} excludes
all other off-diagonal supports, and
Theorem~\ref{thm:allrank} assigns the surviving oriented shared-link element
$t_3=5/612$.  Translation and cubic symmetry make the remaining diagonal a
scalar.  The coefficient is finite range, so the entrywise identity also
holds in every $\mathcal A_\mu$.  This proves
\eqref{eq:g18-exact-infinite-u2}.
\end{proof}

Define
\begin{equation}
 s_{\le2}(u):=\frac83+u+\sigma_2u^2,\qquad t_3:=\frac5{612}.
 \label{eq:g18-s2-definition}
\end{equation}

\begin{theorem}[Exact fixed-spacing carrier sheet on a nonzero-momentum patch]
\label{thm:g18-exact-fixed-k-carrier}
Assume the repaired CMP(1) lemma and
Theorem~\ref{thm:g18-fixed-spacing-band}.  Let $U$ be a relatively compact
open momentum ball whose closure lies in
$\mathbb T^3\setminus\{0\}$, and put
\[
 q_*:=\inf_{k\in U}q(k)>0.
\]
Then there is $u_U>0$ such that for every real
$0<|u|<u_U$ the exact infinite-volume target band over $U$ has a rank-one
analytic Riesz subbundle, descended from $\ker B(k)^\dagger$, separated from
the other two target branches by
\begin{equation}
 \operatorname{gap}_{\rm int}(k,u)
 \ge\frac12t_3u^2q_*,
 \qquad k\in U.
 \label{eq:g18-closed-internal-gap}
\end{equation}
Its direct integral is a reducing, translation-invariant subspace unitarily
equivalent to $L^2(U)$, on which translations and the Hamiltonian act as
\begin{equation}
 (T_xf)(k)=e^{ik\cdot x}f(k),
 \qquad
 (H_{\rm car}(u)f)(k)=E_{\rm car}(k,u)f(k).
 \label{eq:g18-carrier-representation-closed}
\end{equation}
The function $E_{\rm car}(k,u)$ is real analytic in $k$ on $U$.
\end{theorem}

\begin{proof}
Choose $0<r<\rho_0$ in Lemma~\ref{lem:g18-u-analytic-symbol} and set
\[
 M_r:=\sup_{|\zeta|=r}\|h_\zeta\|_{\mu,\sharp}<\infty,
 \qquad C_r:=\frac{2M_r}{r^3}.
\]
Cauchy's estimate in $\mathcal A_\mu$, followed by summing the Taylor tail,
gives for $|u|\le r/2$
\begin{equation}
 h_u(k)=s_{\le2}(u)I_3+t_3u^2B(k)B(k)^\dagger+R_u(k),
 \qquad
 \sup_{k\in\mathbb T^3}\|R_u(k)\|
 \le C_r|u|^3.
 \label{eq:g18-uniform-cubic-remainder}
\end{equation}
Indeed, the $n$th Banach-valued Taylor coefficient has norm at most
$M_rr^{-n}$, and
$\sum_{n\ge3}(|u|/r)^n\le2|u|^3/r^3$ on this disc.

One may therefore take, in addition to the smallness thresholds required by
CMP(1)--CMP(3),
\begin{equation}
 u_U:=\min\left\{\frac r2,
          \frac{t_3q_*}{4C_r}\right\}>0.
 \label{eq:g18-explicit-uU}
\end{equation}
For $0<|u|<u_U$,
\[
 \sup_{k\in U}\|R_u(k)\|
 \le\frac14t_3u^2q_*.
\]
The comparison matrix in
\eqref{eq:g18-uniform-cubic-remainder} has a simple carrier eigenvalue
$s_{\le2}(u)$ and a double eigenvalue
$s_{\le2}(u)+t_3u^2q(k)$.  Weyl's inequality separates the corresponding
exact spectral clusters by at least
\[
 t_3u^2q(k)-2\|R_u(k)\|
 \ge\frac12t_3u^2q_*.
\]
The small contour surrounding the lower cluster therefore defines a rank-one
Riesz projection analytic in $k$.  Put
$P_0(k)=v_0(k)v_0(k)^\dagger$ with
$v_0(k)=w(k)/\sqrt{q(k)}$; this is the limiting comparison projection chosen
by the $BB^\dagger$ coefficient, not a spectral projection of the triply
degenerate $u=0$ fiber.  The same Riesz estimate and
$\|R_u\|/(t_3u^2q_*)=O(|u|/q_*)$ show that, after decreasing $u_U$ if
necessary,
$\sup_{k\in U}\|P_{\rm car}(k,u)-P_0(k)\|<1/2$.  Hence
\[
 v_u(k):=\frac{P_{\rm car}(k,u)v_0(k)}
 {\|P_{\rm car}(k,u)v_0(k)\|}
\]
is a globally defined real-analytic unit section on the whole momentum ball.
Fourier direct integration gives
\eqref{eq:g18-carrier-representation-closed}; analyticity of the simple
eigenvalue follows from analyticity of the matrix symbol.
\end{proof}

\begin{corollary}[No dark carrier on a fixed momentum patch]
\label{cor:g18-no-dark-fixed-patch}
After possibly decreasing $u_U$, the rank-one sheet in
Theorem~\ref{thm:g18-exact-fixed-k-carrier} is seen uniformly by a dressed
cube-boundary Wilson source.  More precisely, put
\[
 v_0(k):=\frac{w(k)}{\sqrt{q(k)}},\qquad k\in U,
\]
and let $S_W(k,u)$ be the normalized literal-Wilson projected synthesis map
written in the orthonormal band frame, with the fixed CMP(4) normalization
chosen so that $S_W(k,0)=I_3$.  If $P_{\rm car}(k,u)$ is the exact rank-one
projection, then $u_U$ may be chosen so that
\begin{equation}
 \inf_{k\in U}
 \bigl\|P_{\rm car}(k,u)S_W(k,u)v_0(k)\bigr\|^2
 \ge\frac14,
 \qquad 0<|u|<u_U.
 \label{eq:g18-fixed-patch-carrier-overlap}
\end{equation}
The number $1/4$ is an absolute projected-norm bound for this normalized
coefficient vector, not a spectral fraction normalized by the full raw-source
covariance.  Consequently every $f\in C_c^\infty(U)$ gives a rapidly decaying
spatial smearing of the literal Wilson plaquettes, with cube-boundary
incidence, whose projection onto the exact carrier sheet is nonzero and
volume independent.
\end{corollary}

\begin{proof}
At $u=0$ the three target levels coincide, so there is no spectrally defined
rank-one carrier projection.  Instead let
$P_0(k)=v_0(k)v_0(k)^\dagger$ be the limiting comparison projection selected
by the $BB^\dagger$ coefficient.  Then
$P_0(k)S_W(k,0)v_0(k)=v_0(k)$.  The remainder estimate
\eqref{eq:g18-uniform-cubic-remainder} and the order-$u^2$ gap imply, by the
finite-dimensional Riesz-projection estimate, that
\[
 \sup_{k\in U}\|P_{\rm car}(k,u)-P_0(k)\|\longrightarrow0
 \quad (u\to0).
\]
CMP(3)--CMP(4) make this explicit.  In the notation
$W=T_cG^{-1/2}$ and $T_W=T_c(I+E)$,
\[
 S_W=W^*T_W=G^{1/2}(I+E),
\]
so $\|S_W-I\|_{\mu,\sharp}\to0$.  Therefore
\[
 \|P_{\rm car}S_Wv_0\|
 \ge1-\|P_{\rm car}-P_0\|-\|S_W-I\|.
\]
Since $\overline U$ is compact, the right-hand side is at least $1/2$ after
shrinking $u_U$.
Squaring gives \eqref{eq:g18-fixed-patch-carrier-overlap}.  Fourier transform
then turns $f\in C_c^\infty(U)$ into a rapidly decaying spatial smearing.
The identity
$f(k)v_0(k)=w(k)f(k)/\sqrt{q(k)}$ identifies its coefficients as the
cube-boundary incidence combination of the three literal Wilson plaquettes.
\end{proof}

\begin{remark}[What has and has not been closed]
\label{rem:g18-internal-sheet-firewall}
Theorem~\ref{thm:g18-exact-fixed-k-carrier} removes the two hypotheses in the
former conditional fixed-$k$ statement: it proves both the exact
orthonormal-symbol coefficient
$[u^2]h_u=\sigma_2I+(5/612)BB^\dagger$ and a remainder uniform in momentum
on a nonzero-momentum patch.  The threshold $u_U$ is
existential because the Yarotsky contraction majorants are existential, and
it necessarily deteriorates as $q_*\downarrow0$.  No rank-one separation at
$\Gamma$ is claimed; the cubic $T_1$ triplet meets there.

For every $a_t>0$, the exact Hamiltonian semigroup restricts to this subspace
as multiplication by $e^{-a_tE_{\rm car}(k,u)}$.  This does \emph{not}
identify $e^{-a_tH_\infty(u)}$ with the finite-temporal-spacing isotropic
Wilson-action transfer matrix.  The anisotropic-time/operator equivalence and
every continuum-limit statement remain separate open problems.
\end{remark}
